As of 2026, the Thesis Center requires an abstract of NO MORE than 250 words. This can be quite limiting, but I would encourage you to write this early this and your title is *required* before anything on your thesis center profile can be saved. There does not seem to be a clear limit on the abstract length in the actual text.

Nulla aliquam sem non dictum pellentesque. Aenean tellus metus, luctus ac nibh eu, luctus fermentum velit. Aliquam id tempus urna. Curabitur non eros non felis euismod mollis vitae a felis. Donec posuere eu velit sed efficitur. Donec lobortis eros velit. Aliquam consectetur convallis est, placerat faucibus velit congue quis. Nulla luctus elit ac lacinia pulvinar. Phasellus tortor libero, tincidunt ut augue sed, hendrerit ultrices sapien. Mauris quam turpis, condimentum in imperdiet vel, iaculis quis dui. Fusce a tortor id eros hendrerit mollis at nec sem.

Cras elementum feugiat diam, accumsan luctus arcu. Suspendisse eleifend lorem ac nunc rhoncus, in lobortis nulla bibendum. Aliquam commodo consequat tempus. Integer ante lacus, tempus hendrerit tortor id, tincidunt semper sapien. Nunc iaculis consequat interdum. Fusce sed elementum velit, eu ornare nunc. Ut luctus tempor pellentesque. Aenean posuere nisl sed dictum rhoncus. Nam leo lorem, scelerisque vel nibh non, condimentum commodo tellus. Morbi consectetur sagittis magna, id vulputate eros venenatis eu. 

% Reducing preventable maternal and child mortality requires more than expanding nominal access to healthcare. In many low- and middle-income countries, populations may live near facilities that lack medicines or diagnostics; services may be formally present but financially difficult to use; and care may be delivered in ways that do not produce client confidence or satisfaction. This dissertation examines how these supply-side conditions—the health service environment—shape maternal and child health service experience and utilization.

% The studies are organized around a common premise: health systems must be evaluated not only by whether care is reached, but by whether the services within reach are prepared, affordable, available, and acceptable. In the first study, I use Service Provision Assessment data from five countries to evaluate the relationship between facility readiness and client satisfaction for antenatal and sick-child care. Disaggregating the Service Readiness Index into its five domains, I find that essential medicines availability was the most consistently associated domain across countries and service types, while others (such as diagnostic capacity and infection prevention) varied in direction and magnitude across settings. Composite readiness scores obscure these distinctions, and domain-specific monitoring offers more actionable guidance. In the second study, I examine whether facilities that charge client fees differ in structural readiness across six countries. Fee charging is most consistently associated with replenishment-dependent domains requiring recurrent procurement—particularly diagnostic capacity and essential medicines—underscoring the need for fee-reduction policies to protect the operating resources that sustain care quality. In the third study, I focus on Malawi and link household and facility data using a probabilistic approach that accounts for the geographic displacement of DHS cluster coordinates. Greater distance and fee exposure are associated with lower use of formal services and, for childhood illness, greater reliance on informal providers; essential medicines readiness shows a positive association with both antenatal and facility-based sick-child care.

% These studies suggest that supply-side conditions are most informative when examined at the domain level, within specific country contexts, and for specific health concerns — aggregate measures consistently obscure more than they reveal. Across this variation, essential medicines availability is the most consistent signal, associated with satisfaction, structural readiness, and care-seeking alike. These findings point toward health system strengthening strategies that tailored for the context and specific health concern, and attentive to the recurrent inputs, that determine whether care within reach is care that works.